\subsection{Zariski's Main Theorem: original maps and separate editorial consequences}\label{algebra-proof-032-zariskis-main-theorem-original-maps-and-separate-editorial-consequences}

Authority: Stacks Project authors, commit a044\allowbreak{}46e5\allowbreak{}7ec1\allowbreak{}fbc2\allowbreak{}52a8\allowbreak{}71af\allowbreak{}cec7\allowbreak{}752f\allowbreak{}b280\allowbreak{}7b14, algebra.\AlgebraIdentifierBreak{}tex:\AlgebraIdentifierBreak{}30454--\AlgebraIdentifierBreak{}30953, SHA-256 FA8B\allowbreak{}B92E\allowbreak{}58A4\allowbreak{}F78A\allowbreak{}2BD0\allowbreak{}1B3B\allowbreak{}6A4A\allowbreak{}87DE\allowbreak{}0A0D\allowbreak{}279F\allowbreak{}5DD9\allowbreak{}0641\allowbreak{}B574\allowbreak{}DD5F\allowbreak{}BFFF\allowbreak{}A4F3. All 23 reports are read. The complete current interval is compared with the authority and the eleven operations in MC-\AlgebraIdentifierBreak{}STK-\AlgebraIdentifierBreak{}ERR-\AlgebraIdentifierBreak{}0649 through 0658. In particular, 0654 adds two references to integrality transitivity. The original human citation keys Henselian and Peskine and the acknowledgement of Thierry Coquand remain intact. Those human sources have not separately been read for this batch.

The source statements, hypotheses and exposition remain identifiable. Every complete argument and strengthening below is separate editorial evidence, not replacement translation text. Three proposed operations add a period and close two occurrences of ``sub algebra''. One contextual sentence-fragment proposal is optional. An omitted argument is not classified as a translation mistake. No novelty claim is made.

\subsubsection{Integral equations with the original localization kernels}\label{algebra-proof-032-integral-equations-with-the-original-localization-kernels}

At 30481--30496 let \(\varphi:R\to S\) and retain the original equation \[
 \sum_{i=0}^{n}\varphi(a_i)t^i=0.
\] If \(n=0\), then \(\varphi(a_0)=0\), so \(\varphi(a_0)t=0\) is integral over \(R\). This is the existing correction 0649; it avoids an unauthorized negative exponent. For \(n\geq1\), put \(z=\varphi(a_n)t\). Multiplying the original equation by \(\varphi(a_n)^{n-1}\) gives exactly \[
 z^n+\sum_{i=0}^{n-1}\varphi(a_i a_n^{\,n-1-i})z^i=0.
\] All exponents are nonnegative, and the equation is monic in \(z\), even when a displayed leading coefficient vanishes or is a zero divisor. Correction 0650 supplies ``by''.

For the trick at 30503--30538 retain the original element \(t_0\), the monic polynomial \(p\), and \(t_0\varphi(p)=\varphi(r)\). Divide \(r=qp+r'\) with \(\deg r'<\deg p\), and define the actual element \[
 \tau=t_0-\varphi(q),\qquad \tau\varphi(p)=\varphi(r').
\] Substitution in the original monic equation for \(t_0\) proves that \(\tau\) is integral over \(R[x]\). If \(p=1\), then \(r'=0\) and \(\tau=0\). Otherwise, in \(S_\tau\) the polynomial \(p(X)-\tau^{-1}r'(X)\) is monic and annihilates the image of \(\varphi(x)\). Thus that image is integral over the subring generated by the image of \(\varphi(R)\) and \(\tau^{-1}\). The image of \(\tau\) is integral over this subring with \(\varphi(x)\) adjoined, by its original integral equation. Transitivity therefore gives, in \(S_\tau\), \[
 \tau^d+\sum_{i=0}^{d-1}\sum_{j=0}^{n_i}
               \varphi(r_{ij})\tau^{i-j}=0,\qquad r_{ij}\in R,\ d\geq1.
\] Choose \(N_0\geq\max(0,j-i)\) over this finite list. The resulting element \[
 E=\tau^{d+N_0}+
       \sum_{i=0}^{d-1}\sum_{j=0}^{n_i}
             \varphi(r_{ij})\tau^{i+N_0-j}\in S
\] maps to zero in \(S_\tau\). By the definition of localization there is \(M\geq0\) with \(\tau^M E=0\) in \(S\). For \(N=N_0+M\) this is the monic equation \[
 \tau^{d+N}+\sum_{i=0}^{d-1}\sum_{j=0}^{n_i}
                \varphi(r_{ij})\tau^{i+N-j}=0
\] in the original ring. Every other exponent is smaller than \(d+N\). No injectivity of \(S\to S_\tau\), reducedness or cancellation was assumed.

For 30541--30568 retain \(a=\varphi(a_k)\). If \(a\) is nilpotent in \(S\), choose \(n\) with \(a^n=0\) and \(q=0\); the required element is zero. Otherwise localize \(R\) at \(a_k\) and \(S\) at \(a\), retaining the actual maps \(\varphi'\) and \(t'\). The polynomial \(p/a_k\) is monic over \(R_{a_k}\), and the preceding argument gives \(q'\) such that \(t'-\varphi'(q')\) is integral. Choose \(n,q\) with \(a_k^n q'\) equal to the image of \(q\). Then \[
 b=a^n t-\varphi(q)\in S
\] has integral image in \(S_a\). Here is the denominator and kernel calculation behind the cited localization lemma. Write its integral equation as \[
 b^d+\sum_{i=0}^{d-1}\frac{\varphi(r_i)}{a^{e_i}}b^i=0
 \quad\text{in }S_a.
\] Choose \(N\) such that \(N(d-i)\geq e_i\) for all \(i\). For \(y=a^N b\) the element \[
 E=y^d+\sum_{i=0}^{d-1}\varphi(r_i)a^{N(d-i)-e_i}y^i
\] maps to zero. Choose \(M\geq0\) with \(a^M E=0\). Multiplication by \(a^{Md}\), which also kills \(E\), gives for \(z=a^M y\) \[
 z^d+\sum_{i=0}^{d-1}\varphi(r_i)a^{(N+M)(d-i)-e_i}z^i=0.
\] Consequently the actual element \[
 z=a^{N+M+n}t-\varphi(a_k^{N+M}q)
\] is integral over \(R\). This retains both the polynomial denominator and the possible localization kernel. The original localization result at 7596--7627 was reread; its original multipliers remain valid.

\subsubsection{The conductor and the coefficient argument}\label{algebra-proof-032-the-conductor-and-the-coefficient-argument}

In the original situation at 30571 set \(A=\operatorname{Im}(\varphi:R[x]\to S)\). The conductor is \[
 J=\{g\in S:gS\subseteq A\}.
\] Because \(1\in S\), every member lies in \(A\). Addition and multiplication by \(S\) preserve its defining condition, so it is an ideal of \(S\) and of \(A\), and \[
 J=\operatorname{Ann}_A(S/A).
\] The module \(S/A\) is finite over \(A\). For any specified \(g\in A\), \[
 A_g=S_g\quad\Longleftrightarrow\quad
                  g^N\in J\text{ for some }N\geq0.
\] Indeed localization of \(0\to A\to S\to S/A\to0\) is exact. If its last term vanishes, each member of a finite generating list is killed by a power of \(g\); their maximum kills the whole module. The converse follows by inverting \(g\). This is a statement about the whole principal localization, not about an individual pair of local rings at a chosen prime.

The ``In Situation \ldots'' phrase at 30584 is a conventional contextual fragment; changing its punctuation is optional. Correction 0651 already states correctly that \(\varphi(R)\) is integrally closed in \(S\). For the leading-coefficient lemma, choose a finite \(A\)-module generating list \(t_1,\ldots,t_\ell\) for \(S\), adjoining \(1\) if necessary so that the list is nonempty even for the zero ring. Every \(ut_i\) is integral over \(R[x]\). If \(u\varphi(P)\in J\), then \[
 \varphi(P)ut_i\in A.
\] The preceding lemma gives \(n_i,q_i\) such that \[
 \varphi(a_k^{n_i})ut_i-\varphi(q_i)
\] is integral over \(R\), hence belongs to \(\varphi(R)\). Thus \(\varphi(a_k^{n_i})ut_i\in A\), and \(m=\max_i n_i\) gives \(u\varphi(a_k)^m\in J\). Multiplication by the remaining powers takes place in \(A\); no division occurs.

For 30607--30623 suppose \(u\varphi(P)\in\sqrt J\). Choose the actual \(n\geq1\) with \(u^n\varphi(P)^n\in J\). The leading coefficient of \(P^n\) is \(a_k^n\), so the previous argument gives \[
 u^n\varphi(a_k)^{nm}\in J
\] for some \(m\geq1\), increasing \(m\) if needed. Multiplying by \(u^{n(m-1)}\) gives \((u\varphi(a_k))^{nm}\in J\). Subtract \(u\varphi(a_k)x^k\) from \(u\varphi(P)\) and induct on \(k\); for \(k=0\) the assertion is the hypothesis. Each original coefficient therefore satisfies \(u\varphi(a_i)\in\sqrt J\). This uses powers of specified elements, not a claim that the entire radical ideal has one nilpotence exponent.

\subsubsection{Strong transcendence and the reducedness boundary}\label{algebra-proof-032-strong-transcendence-and-the-reducedness-boundary}

The source definition at 30629 is retained: for an inclusion \(R\subseteq S\), \(x\in S\) is strongly transcendental if \[
 uP(x)=0,\quad P(T)=\sum_{i=0}^k a_iT^i\in R[T]
 \quad\Longrightarrow\quad ua_i=0\text{ for every }i.
\] Equivalently, for every such original polynomial, \[
 \operatorname{Ann}_S(P(x))
       =\bigcap_{i=0}^k\operatorname{Ann}_S(a_i).
 \tag{ST}
\] The right side is always contained in the left. The reverse is exactly the definition. These are the kernels of the two actual \(S\)-linear maps \[
 S\longrightarrow S,\ u\longmapsto uP(x),
 \qquad
 S\longrightarrow S^{k+1},\ u\longmapsto(ua_0,\ldots,ua_k).
\] Tensoring their kernel sequences by a flat \(S\)-algebra \(T\) shows that (ST) is preserved for the image of \(x\) over the image of the same original \(R\). Every coefficient in that image has an \(R\)-lift, so this verifies all its polynomials. If \(T\) is faithfully flat, equality of these kernels after tensoring forces equality before tensoring: the quotient of the first kernel by the second has zero faithful tensor. Thus strong transcendence is reflected as well. These assertions do not claim invariance under enlarging the coefficient ring by an arbitrary \(R\)-base change.

In particular localization at any multiplicative subset \(W\subset S\) preserves strong transcendence over the image of \(R\). Directly, if \((u/w)P(x)=0\), some \(w'\in W\) kills \(uP(x)\) in \(S\); hence \(w'ua_i=0\) for every \(i\), so every localized coefficient product vanishes.

\textbf{Editorial strengthening of 30644--30663, ALGEBRA-\AlgebraIdentifierBreak{}RECON-\AlgebraIdentifierBreak{}762.} Reducedness is unnecessary for passage to a minimal prime. Let \(\mathfrak q\) be any minimal prime of \(S\), and put \(\mathfrak p=R\cap\mathfrak q\). Suppose \[
 uP(x)\in\mathfrak q.
\] If \(u\in\mathfrak q\), every \(ua_i\) already belongs to \(\mathfrak q\). Otherwise \(u\) is a unit in \(S_{\mathfrak q}\), so \(P(x)\) belongs to \(\mathfrak qS_{\mathfrak q}\). This local ring has exactly one prime, because primes of the localization correspond to primes contained in the minimal prime \(\mathfrak q\). Its nilradical is therefore \(\mathfrak qS_{\mathfrak q}\). Choose \(N\geq1\) with \(P(x)^N=0\) there.

Let \(R_0\) be the image of \(R\) in \(S_{\mathfrak q}\), retaining its possible nilpotents. Localized strong transcendence applied with multiplier \(1\) to the polynomial \(P(T)^N\) says that all its coefficients are zero in \(R_0\). Hence \(P(T)\) is nilpotent in \(R_0[T]\). For every prime \(\mathfrak r\) of \(R_0\), its image in the polynomial ring over the domain \(R_0/\mathfrak r\) must be zero; a polynomial ring over a domain has no nonzero nilpotents, since leading coefficients of products do not vanish. Thus every coefficient of \(P\) lies in every \(\mathfrak r\), and is nilpotent in \(R_0\). Its image in \(S_{\mathfrak q}\) is nilpotent and belongs to \(\mathfrak qS_{\mathfrak q}\). Contraction gives \(a_i\in\mathfrak q\), and hence \(ua_i\in\mathfrak q\). This proves strong transcendence in \(S/\mathfrak q\) over \(R/\mathfrak p\), without discarding any nilpotent term in the argument.

The reduced source proof remains valid and shorter: \(S_{\mathfrak q}\) is then a field, and the original multiplier \(u'\notin\mathfrak q\) kills \(uP(x)\); strong transcendence yields \(uu'a_i=0\), and primality yields \(ua_i\in\mathfrak q\). The general argument is not substituted into that translation.

When \(S\) is reduced, this also gives the exact componentwise characterization: \[
 x\text{ strongly transcendental over }R
 \quad\Longleftrightarrow\quad
 x\bmod\mathfrak q\text{ transcendental over }
 R/(R\cap\mathfrak q)
 \text{ for every minimal }\mathfrak q.
\] Here transcendence over the domain is equivalent to transcendence over its fraction field, after clearing denominators. Forward implication follows from the theorem. Conversely \(uP(x)=0\) implies \(ua_i\in\mathfrak q\) for every minimal prime, by the domain criterion, and their intersection is the nilradical, zero by reducedness.

Reducedness cannot be dropped from this converse. Take \[
 R=k,\qquad S=k[X,\epsilon]/(\epsilon^2,X\epsilon),\qquad x=X.
\] As a \(k\)-module this is \(k[X]\oplus k\epsilon\), with multiplication \[
 (f,a\epsilon)(g,b\epsilon)
  =(fg,(f(0)b+g(0)a)\epsilon).
\] The direct expression proves \(\epsilon\ne0\), \(\epsilon^2=0\) and \(X\epsilon=0\). Its unique minimal prime is \((\epsilon)\), because the quotient is \(k[X]\), and every prime contains the nilpotent \(\epsilon\). The image of \(X\) in that quotient is transcendental. Nevertheless \(\epsilon X=0\) violates strong transcendence for \(P(T)=T\), whose coefficient \(1\) is not annihilated by \(\epsilon\). Arbitrary prime quotients also need not preserve strong transcendence: \(X\) is strongly transcendental in the domain \(k[X]\) over \(k\), but its image in \(k[X]/(X)\) is zero.

\subsubsection{The original nowhere quasi-finite argument}\label{algebra-proof-032-the-original-nowhere-quasi-finite-argument}

At 30666--30711 suppose \(R\subset S\) are domains, \(x\) is transcendental over \(R\), and \(S\) is finite over the actual embedded polynomial ring \(R[x]\). First assume \(R\) normal. The original polynomial-normality lemma at 8246--8257 makes \(R[x]\) normal. If \(R\to S\) were quasi-finite at \(\mathfrak q\), its residue extension over \(\mathfrak p=R\cap\mathfrak q\) would be finite. For \(\mathfrak r=R[x]\cap\mathfrak q\), the intermediate residue extension is also finite, so \[
 \mathfrak pR[x]\subsetneq\mathfrak r.
\] Equality would give the transcendental residue field \(\kappa(\mathfrak p)(x)\). Going down for the integral inclusion \(R[x]\subset S\) produces \(\mathfrak q_0\subsetneq\mathfrak q\) over \(\mathfrak pR[x]\). Both primes lie in the same \(R\)-fibre. Every open containing \(\mathfrak q\) contains its generization \(\mathfrak q_0\); thus \(\mathfrak q\) is not isolated in that fibre, contradicting quasi-finiteness.

The going-down statement at 8517--8572 was reread. Its previously recorded proof corrections in ALGEBRA-\AlgebraIdentifierBreak{}RECON-\AlgebraIdentifierBreak{}216 through 218 retain the minus sign \(g^n=-\sum_{j<n}a_jg^j\) and handle \(z=0\) before scaling the minimal polynomial. Those proof corrections are dependencies, not newly allocated operations or instructions to expand a translation.

For arbitrary \(R\), retain \(K=\operatorname{Frac}(R)\), \(L=\operatorname{Frac}(S)\), the integral closure \(\overline R\subset K\), and \(\overline S=S[\overline R]\subset L\). The actual multiplication map \[
 S\otimes_R\overline R\longrightarrow\overline S,\qquad
 s\otimes r\longmapsto sr
\] is surjective. Hence \(\overline S\) is finite over \(\overline R[x]\), by transporting the original finite module generators. Transcendence is preserved because \(\operatorname{Frac}(\overline R)=K\). The inclusion \(S\subset\overline S\) is integral: any specified element lies in an algebra generated by finitely many elements integral over \(R\), hence integral over \(S\). Lying over provides a prime above every \(\mathfrak q\). The original four-ring lemma would transfer quasi-finiteness to that prime, contradicting the normal case. No finiteness of the whole normalization is assumed.

For the source reduced-ring lemma at 30714--30732, choose a minimal prime \(\mathfrak q_0\subseteq\mathfrak q\), pass to the domain inclusion \[
 R/(R\cap\mathfrak q_0)\ \subseteq\ S/\mathfrak q_0,
\] and use the preceding domain theorem and the four-ring lemma. Correction 0652 fixes only ``be'' to ``is''.

The general minimal-prime result in ALGEBRA-\AlgebraIdentifierBreak{}RECON-\AlgebraIdentifierBreak{}762 proves the same nowhere quasi-finite conclusion for any inclusion \(R\subseteq S\), with no reducedness assumption, provided \(x\) is strongly transcendental and \(S\) is finite over \(R[x]\). The finite module generators descend to the displayed domain quotient, and the quotient comparison map in the four-ring lemma is surjective. Thus hypothetical quasi-finiteness at \(\mathfrak q\) would give the forbidden quasi-finiteness at its quotient prime. This is the precise propagation of the strengthening to the next source lemma; the original reduced hypotheses remain in the translated statements.

\subsubsection{Monogenic localization and the explicit Horner witness}\label{algebra-proof-032-monogenic-localization-and-the-explicit-horner-witness}

At 30735--30776 let \(S=R[x]/I\), let \(B\) denote the source integral closure \(S'\) of \(R\) in \(S\), and suppose \(R\to S\) is quasi-finite at \(\mathfrak q\). There must be an original \(f=\sum a_iX^i\in I\) with some \(a_i\notin\mathfrak p=R\cap\mathfrak q\): otherwise the fibre ideal is zero and its ring is \(\kappa(\mathfrak p)[X]\), which the domain theorem shows has no quasi-finite point over \(\kappa(\mathfrak p)\). Its coefficients give a relation \[
 \sum_{i=0}^m b_i x^i=0,\qquad b_i\in B,
 \qquad b_j\notin\mathfrak q\text{ for some }j.
\] When \(m=0\) this is impossible; correction 0653 removes the duplicated copula, without altering that argument.

For \(m\geq1\) the first integral-equation lemma shows that \(b_mx\) is integral over \(B\). Since \(B\) is integral over \(R\), transitivity makes \(b_mx\) integral over \(R\), and its membership in \(S\) then puts it in \(B\). These are precisely the two transitivity insertions already composed as 0654. If \(b_m\notin\mathfrak q\), the inclusion \(B_{b_m}\subset S_{b_m}\) is surjective because \(x=(b_mx)/b_m\); it is injective by exactness of localization of the original inclusion. Otherwise set \(b'_{m-1}=b_mx+b_{m-1}\in B\), and retain the relation \[
 b'_{m-1}x^{m-1}+b_{m-2}x^{m-2}+\cdots+b_0=0.
\] The replacement leading coefficient has the same residue as \(b_{m-1}\) at \(\mathfrak q\). At least one coefficient is still outside that prime, and induction applies. The proposed period at 30770 is punctuation only.

\textbf{Editorial consequence ALGEBRA-\AlgebraIdentifierBreak{}RECON-\AlgebraIdentifierBreak{}763.} The source induction produces a specified witness from the original relation. Put \[
 j=\max\{i:b_i\notin\mathfrak q\}.
\] One has \(j\geq1\), since otherwise reduction of the original relation modulo \(\mathfrak q\) would assert \(b_0=0\) there. Define \(c_m=b_m\) and, successively for \(r=m,m-1,\ldots,j+1\), \[
 c_{r-1}=c_rx+b_{r-1}.
\] At each stage the retained relation is \[
 c_rx^r+\sum_{i=0}^{r-1}b_ix^i=0.
\] The integral-equation lemma and transitivity prove \(c_rx\in B\), and hence \(c_{r-1}\in B\). For \(r>j\), induction also gives \(c_r\in\mathfrak q\). Therefore the final elements are \[
 g=c_j=\sum_{i=j}^m b_i x^{i-j}\in B\setminus\mathfrak q,
 \qquad h=gx\in B.
\] The assertion about \(h\) uses the last relation of degree \(j\geq1\). Every summand and the original coordinate \(x\) are retained.

The actual finite \(R\)-subalgebra \(A=\operatorname{Im}(R)[g,h]\subset S\) satisfies \(A_g=S_g\): its localization contains \(x=h/g\), and it already contains the image of \(R\). The canonical localization inclusion is injective, giving the inverse to this generator description. If specified monic equations for \(g,h\) have degrees \(r,s\), successive substitution of their leading powers expresses every element of \(A\) as an \(R\)-linear combination of \[
 g^a h^b,\qquad 0\leq a<r,\quad 0\leq b<s.
\] Thus \(rs\) is a bound on the length of this specified generating list; linear independence is not asserted.

\subsubsection{The main induction with both final denominators}\label{algebra-proof-032-the-main-induction-with-both-final-denominators}

The theorem at 30779 retains \(S\) as the finite-type \(R\)-algebra, as already corrected by 0655. Corrections 0656--0658 retain ``theorem'', the example comma and the localization explanation. Adding a verb to the elliptical example is optional; it is not a remaining mathematical defect. The two proposed spelling changes occur at 30791 and 30857; the differently arranged phrase at 30894 is not automatically changed.

Write \(B\subset S\) for the original integral closure of \(R\), and choose the original \(x_1,\ldots,x_n\) with \(S\) finite over \(R[x_1,\ldots,x_n]\). If \(n=0\), then \(B=S\), and \(g=1\) works.

For \(n=1\), the permanence lemma allows the actual intermediate base inclusion \(B\subset S\), and the same module generators show finiteness over \(B[x_1]\). Prove the assertion with this base, denoting it by \(R\) in the source comparison: \(R\subset S\) is now integrally closed. This is an explicit intermediate-ring map, not an identification of the old base ring with its integral closure. Let \(A=\operatorname{Im}(R[x]\to S)\) and retain its conductor \(J\). The coefficient lemma proves that the image of \(x\) in \[
 S/\sqrt J
\] is strongly transcendental over the embedded ring \(R/(R\cap\sqrt J)\). These two rings are reduced. The original reduced nowhere quasi-finite lemma therefore suffices; the stronger editorial version is not needed to replace this proof. The four-ring lemma implies \(\mathfrak q\notin V(J)\).

Choose the actual \(s\in J\setminus\mathfrak q\), \(s=\varphi(f)\in A\). The inclusion \(A_s\to S_s\) is an isomorphism: any \(y/s^n\) is the image of \((sy)/s^{n+1}\), whose numerator is in \(A\), and localization preserves the inclusion. This explicit formula is independent of the representation because it is an inverse to the canonical inclusion. The local fibre criterion transfers quasi-finiteness to the corresponding prime of \(A\). The monogenic lemma, using that \(R\) is integrally closed in \(A\), gives \(h\in R\setminus\mathfrak q\) with \(R_h=A_h\).

Write the image of the original \(s\) in \(A_h=R_h\) as \(a/h^N\), \(a\in R\). Since \(s\) and \(h\) avoid \(\mathfrak q\), so does \(a\). Set \(g=ha\). In any localization where \(g\) is invertible, both \(h,a\) are invertible and the equality \(s=a/h^N\) makes \(s\) invertible. Conversely, after inverting \(h,s\), the element \(a=s h^N\) and hence \(g\) are invertible. The universal properties of these localizations give the canonical equalities \[
 S_g=S_{hs}=A_{hs}=R_{ha}=R_g.
\] All equalities mean the maps induced by the original inclusions and these stated inverse fractions. No unsupported cancellation in the original rings occurs.

For \(n>1\), retain the source \(R'\), the integral closure of \(R[x_1,\ldots,x_{n-1}]\) in \(S\). Permanence transfers quasi-finiteness, and the same generators make \(S\) finite over \(R'[x_n]\). The case \(n=1\) gives \(g'\in R'\setminus\mathfrak q\) with \(R'_{g'}=S_{g'}\). This integral closure need not be of finite type over \(R\).

Choose a finite \(R\)-algebra generating list for \(S_{g'}\) in the form \(y_i/(g')^{n_i}\), and adjoin \(g'\) to it if needed. Define the actual \(R\)-subalgebra \[
 R''=\operatorname{Im}(R)[x_1,\ldots,x_{n-1},y_1,\ldots,y_N,g']
 \subset R'.
\] Localization of its inclusion is injective; the specified generators prove \(R''_{g'}=S_{g'}\). Each adjoined generator is integral over \(R[x_1,\ldots,x_{n-1}]\), so its monic equation supplies a finite bounded-power spanning list, and their products prove that \(R''\) is finite over that original polynomial-image ring. For \(\mathfrak q''=R''\cap\mathfrak q\), one has \(R''_{\mathfrak q''}=S_{\mathfrak q}\). Since \(R''\) is of finite type over \(R\), the local fibre criterion transfers quasi-finiteness again.

Induction gives \(C=R'''\subset R''\), integral over \(R\), and \(g''\in C\setminus\mathfrak q\) with \(C_{g''}=R''_{g''}\). Retain the source numerator \(g'''\in C\) and exponent \(N\) in \[
 g'=\frac{g'''}{(g'')^N}\quad\text{in }R''_{g''}.
\] The numerator avoids \(\mathfrak q\), by reduction in the localization at that prime. Set \(g=g''g'''\). After inverting \(g\), both \(g''\) and \(g'\) become units. Localizing the two previously proved canonical equalities gives \[
 C_g=R''_g=S_g.
\] Every element of \(C\) is integral over \(R\), so \(C\subseteq B\subseteq S\); localization preserves these inclusions and gives \(B_g=S_g\). This proves the source induction with both final denominators retained.

\subsubsection{Finite envelopes and the exact open immersion}\label{algebra-proof-032-finite-envelopes-and-the-exact-open-immersion}

For the openness lemma at 30880 and the global lemma at 30906, retain the original \(B=S'\). Given a local witness \(B_g=S_g\), choose the original finite \(R\)-algebra generators \(z_1,\ldots,z_r\) of \(S\). Express their images in \(B_g\) as \(a_j/g^{n_j}\), \(a_j\in B\), and set \[
 A=\operatorname{Im}(R)[g,a_1,\ldots,a_r]\subseteq B.
\] Every generator is integral over \(R\). Given monic equations of degrees \(d_0,d_1,\ldots,d_r\), the monomials \[
 g^{e_0}a_1^{e_1}\cdots a_r^{e_r},
 \qquad 0\leq e_i<d_i,
\] span \(A\) over \(R\), by replacing one excessive exponent at a time with its stated monic equation. This is a finite list of size \(\prod_i d_i\), with no freeness claim. The actual inclusion identifies \(A_g=S_g\), since all \(z_j\) and \(g^{-1}\) lie in \(A_g\). Finite \(A/R\) has finite fibres, and localization at \(g\) preserves quasi-finiteness at each of its primes. Thus every point of \(D_S(g)\) is quasi-finite over \(R\), proving openness.

Now assume \(S/R\) is quasi-finite at every prime. Choose finitely many original witnesses \(g_i\in B\) with \[
 B_{g_i}=S_{g_i},\qquad (g_1,\ldots,g_n)S=S.
\] For \(S=0\), also \(B=0\); take \(A=0\), whose spectrum is empty and whose map to \(S\) is the identity. Its finite presentation over itself, flatness and epimorphism assertions below are immediate. For \(S\ne0\), choose \(n\geq1\) as above.

For every \(i\), clear denominators in the same finite list of \(R\)-algebra generators of \(S\), obtaining elements \(a_{ij}\in B\) with \(z_j=a_{ij}/g_i^{n_{ij}}\) in \(S_{g_i}\). Put \[
 A=\operatorname{Im}(R)[g_1,\ldots,g_n,(a_{ij})_{i,j}]\subseteq B.
\] The preceding bounded-power argument proves \(A\) finite over \(R\), and \(A_{g_i}=S_{g_i}\) for every \(i\). No assertion that the \(g_i\) generate the unit ideal in \(A\) is made.

\textbf{Editorial consequence ALGEBRA-\AlgebraIdentifierBreak{}RECON-\AlgebraIdentifierBreak{}764.} The actual map \[
 \operatorname{Spec}(S)\longrightarrow\operatorname{Spec}(A)
\] is an open immersion onto \(U=\bigcup_iD_A(g_i)\). Every source prime avoids some \(g_i\), since their ideal in \(S\) is the unit ideal. On that chart the canonical ring isomorphism \(A_{g_i}=S_{g_i}\) gives a homeomorphism and isomorphisms of all local rings. Conversely every prime in \(D_A(g_i)\) has its unique preimage on this chart. If two source primes had the same image, choosing such a \(g_i\) for the image puts both in this one chart and identifies them. These chart inverses agree on intersections, because their ring maps on \(A_{g_ig_j}=S_{g_ig_j}\) are all induced by the same inclusion. They therefore give the inverse homeomorphism and inverse maps on the structure sheaves. This proves the open immersion, not merely its underlying topological assertion. Replacing \(A\) by \(B\) proves the analogous statement for \(\operatorname{Spec}(B)\).

Here is the full module argument for source assertion (2). If \(D_A(g)\subseteq U\), then the \(g_i\) generate the unit ideal in \(A_g\); otherwise a maximal ideal containing them would be a point of \(D_A(g)\) outside \(U\). Localizing \(A_g\to S_g\) further at each \(g_i\) gives an isomorphism. Its kernel and cokernel vanish after each localization. For any module \(M\) over a ring \(D\) and finite elements \(v_i\) generating the unit ideal, \(M_{v_i}=0\) for all \(i\) implies \(M=0\): for a specified \(m\in M\), choose \(N_i\) with \(v_i^{N_i}m=0\). The powers \(v_i^{N_i}\) still generate the unit ideal, since a maximal ideal containing them would contain all \(v_i\). Taking a linear combination equal to \(1\) proves \(m=0\). Apply this to the kernel and cokernel. Thus \(A_g=S_g\), and the same proof works for \(B\). This completes the source's omitted details, without adding a Noetherian hypothesis. The original cover lemma at 4133--4217 was reread.

Both \(A\to S\) and \(B\to S\) are flat epimorphisms of finite presentation. We prove each property with the original \(g_i\). Write \(D=A\) or \(B\). To prove flatness, tensor any injection of \(D\)-modules by \(S\). Its kernel, viewed as an \(S\)-module, vanishes after localizing at each \(g_i\), because \(S_{g_i}=D_{g_i}\) is a flat localization of \(D\). Since the \(g_i\) generate the unit ideal in \(S\), the preceding module argument makes the kernel zero.

For the epimorphism, consider \[
 \mu:S\otimes_D S\longrightarrow S,\qquad s\otimes t\longmapsto st.
\] The elements \(g_i\otimes1\) generate the unit ideal in its source: transport an actual equality \(1=\sum c_i g_i\) in \(S\) through \(s\mapsto s\otimes1\). After localizing at \(g_i\otimes1\), the source is \[
 S_{g_i}\otimes_D S=D_{g_i}\otimes_D S=S_{g_i},
\] and \(\mu\) is the identity under these canonical maps. Its kernel and cokernel therefore vanish by the same finite-cover argument. The inverse is \(s\mapsto s\otimes1\), and this also proves \(s\otimes1=1\otimes s\). If two ring maps from \(S\) agree on \(D\), applying their induced map from this tensor product to that equality shows they agree on every \(s\); this is the categorical epimorphism assertion.

For finite presentation, choose actual \(D\)-algebra generators \(z_1,\ldots,z_r\) of \(S\); the original \(R\)-generators suffice. In the unit equality \(1=\sum_i c_i g_i\), choose polynomials \(C_i\in D[X_1,\ldots,X_r]\) with \(C_i(z)=c_i\). Put \[
 F(X)=\sum_i g_i C_i(X)-1,\qquad
 Q=D[X_1,\ldots,X_r]/(F).
\] The evaluation map \(Q\to S\) is surjective, and the \(g_i\) generate the unit ideal of \(Q\). For each \(i,j\), write \(z_j=a_{ij}/g_i^{n_{ij}}\) in \(S_{g_i}=D_{g_i}\), retaining the corresponding original numerator and denominator. Choose \(m_{ij}\geq0\) killing the image of \(g_i^{n_{ij}}X_j-a_{ij}\) under evaluation; this is possible because that image vanishes after inverting \(g_i\). Thus \[
 K_{ij}=g_i^{m_{ij}}(g_i^{n_{ij}}X_j-a_{ij})
\] belongs to the actual evaluation kernel \(K\subset Q\). In \(Q_{g_i}\), the finitely many \(K_{ij}\) for this fixed \(i\) generate \(K_{g_i}\): the original polynomial evaluation into \(D_{g_i}\) has kernel generated by \(X_j-a_{ij}/g_i^{n_{ij}}\), as follows by substituting the variables one at a time in each polynomial. The relation \(F\) already evaluates to zero. Let \(K_0\) be the ideal of \(Q\) generated by all \(K_{ij}\). Then \((K/K_0)_{g_i}=0\) for every \(i\), and the finite unit cover in \(Q\) gives \(K=K_0\). We have the explicit finite presentation \[
 S\cong
 D[X_1,\ldots,X_r]\big/
 \left(F,\,
       \bigl(g_i^{m_{ij}}
             (g_i^{n_{ij}}X_j-a_{ij})\bigr)_{i,j}\right).
\] Every localization-kernel power remains present. This proves finite presentation over \(D=A\) or \(B\); it does not assert finite presentation over the original \(R\), since the finite algebra \(A/R\) need not be finitely presented.

Finally the same charts give the exact ring equalizer \[
 S\longrightarrow\prod_iD_{g_i}
       \rightrightarrows\prod_{i,j}D_{g_ig_j},
 \tag{C}
\] where the two arrows restrict the \(i\)- and \(j\)-components, and the first map uses the original inclusions. Injectivity follows from the unit-cover module argument in \(S\). For existence, identify \(D_{g_i}=S_{g_i}\) and let \((s_i)\) be a compatible tuple. Clear the finitely many denominators to choose \(N_0\) and \(b_i^0\in S\) representing \(g_i^{N_0}s_i\). In \(S_{g_j}\), the difference \(b_i^0-g_i^{N_0}s_j\) vanishes after inverting \(g_i\), by compatibility. Choose one \(M\) at least as large as all the finitely many required annihilating exponents, and set \[
 N=N_0+M,\qquad b_i=g_i^M b_i^0.
\] Then \(b_i=g_i^N s_j\) in \(S_{g_j}\) for every \(i,j\). The elements \(g_i^N\) generate the unit ideal in \(S\); choose actual \(d_i\in S\) with \(\sum_i d_i g_i^N=1\). The element \(s=\sum_i d_i b_i\) restricts to \(s_j\) for every \(j\), proving (C) with the original gluing maps.

\subsubsection{Reading, propagation and remaining scope}\label{algebra-proof-032-reading-propagation-and-remaining-scope}

The new corpus query ``Zariski Main conductor'' was inspected in both layers. Its two research records route to a TeX and a PDF record of the Bianchi modular-forms item by Vlad Serban; they are not counted as two independently read works. The local-path layer routes to an original Connes--Consani TeX copy and an Arithmetic Time PDF. All four hits remain unread and unadopted. No PDF fallback or novelty inference occurs. The route preserves the previous seven separately recorded human-source readings and updates only the bounded Stacks coverage.

ALGEBRA-\AlgebraIdentifierBreak{}RECON-\AlgebraIdentifierBreak{}762 propagates its minimal-prime strengthening to the source's next nowhere quasi-finite lemma and records the precise reducedness boundary. ALGEBRA-\AlgebraIdentifierBreak{}RECON-\AlgebraIdentifierBreak{}763 combines the retained integral-equation correction with the original induction, producing the actual Horner numerator and a two-integral-generator finite algebra. ALGEBRA-\AlgebraIdentifierBreak{}RECON-\AlgebraIdentifierBreak{}764 uses the original main theorem and finite-cover arguments to prove the finite envelope, actual open immersion, flat epimorphism, finite presentation and equalizer. Their complete proofs remain here, with zero translation replacement operations.

The reviewed four-ring results ALGEBRA-\AlgebraIdentifierBreak{}RECON-\AlgebraIdentifierBreak{}741, 742 and 747 supply the original comparison and its already recorded finite-map strengthening; only the source surjective version is needed in the main argument. The finite-locus boundary in ALGEBRA-\AlgebraIdentifierBreak{}RECON-\AlgebraIdentifierBreak{}748 remains intact: these open immersions are into the integral closure and finite intermediate algebras, not an assertion that every quasi-finite map is finite near every point of the original base. Nor does this batch claim a global finite envelope for a possibly non-quasi-compact quasi-finite locus of an arbitrary finite-type map. The constructed global envelope is proved under the source's stated hypothesis that all of \(S/R\) is quasi-finite.

Next is Applications of Zariski's Main Theorem, authority 30954. Lookahead 30954--30975 is not adjudicated. Broader synthesis follows the core task. No source or translation mutation, TeX build or publication is performed by this review.
